There are two men in the governor's hall. One can order an execution. The other has been delivered into his hands. If we were looking for the man with power, the answer would seem obvious. Yet today the Church asks us to recognize our King in the prisoner.

Pilate wants to know whether Jesus is a king. It is a practical question. A king might have followers willing to fight for him. He might threaten the authority Pilate represents. Jesus points to something Pilate can see: his servants are not fighting to prevent his being handed over. Then he speaks about truth, and about those who hear his voice.

That is a strange way for a prisoner to answer a question about power. But it brings us to the heart of this feast. Christ reigns by giving his people a new life. He makes them capable of hearing him, and that hearing changes how they live with one another.

Saint Augustine notices a small distinction in Jesus' answer. Jesus says that his kingdom is not of this world. He does not say that it is absent from this world. The distinction matters. The King's people are here. They have homes and work, neighbors and responsibilities. They live among people who believe and people who do not. Christ does not take them out of those relationships. He gives them a different source of life within them.

So when we hear that his kingdom is not of this world, we cannot put it safely away in a distant heaven, leaving our ordinary conduct untouched. The question becomes much closer: what is shaping my life here? Whose voice teaches me what to love, what to refuse, what to do when I am wrong?

The Epistle gives us a way to answer. It begins with thanksgiving. The Father has rescued us from darkness and brought us into the kingdom of his beloved Son. Augustine uses this very passage to explain the Gospel. We become the King's people through a new birth. We owe our belonging to a gift.

Saint John Chrysostom draws attention to how generous that gift is. God gives the inheritance, and he also makes us fit to receive it. Someone can be given an office without having the ability to carry it out. God reaches more deeply. His grace changes the recipient. The King gives us the life from which obedience can grow.

That matters especially when obedience seems beyond us. We may know very well that we ought to stop nursing a grievance, tell the truth about a failure, or receive a correction without immediately striking back. Knowing the right thing and wanting to do it can seem far apart. The Epistle begins with what God has done for people who needed rescue. Our weakness is a reason to turn toward him. The life we need begins in his generosity.

Then listen to the final movement of the Gospel. Jesus connects belonging to truth with hearing his voice. Augustine explains that hearing means inward belief and obedience. Pilate can hear every sound Jesus makes. The hearing Christ asks of us reaches further: it allows his truth to judge our choices and change them.

Augustine also notices which comes first. Belonging to truth gives rise to hearing. We do not earn God's first gift by proving how good we already are at listening. Grace enables the response. That is why Christian obedience can be both humble and active. We really must respond, and we can give thanks for the very possibility of responding.

Consider what happens when someone corrects us. We may listen closely, but only to prepare our defense. We can remember every word and remain determined that none of it will change us. Hearing the King asks something more costly. It asks us to let truth reach the place we are protecting. Perhaps there is a real fault beneath the excuse. Perhaps the person we resent has told us something we needed to hear.

The King who asks this of us has given himself for us. The Epistle identifies him as the one through whom all things were created and in whom they continue to exist. Then it brings that immense claim to the Cross. The Lord of creation makes peace through his own blood. His refusal of an armed rescue belongs to that saving work. The prisoner before Pilate is moving toward the sacrifice through which enemies can be reconciled.

We can therefore trust his correction. He seeks the healing of the person he calls to obedience. The Mass's opening prayer names sin as the wound that divides the human family. Its appeal to Christ's gentle dominion concerns that wound. A divided heart needs more than instructions delivered from a distance. It needs the King whose self-giving can restore it.

At the offering, the prayer called the Secret asks Christ for unity and peace among the nations. The one offered in sacrifice is himself the giver of the peace being sought. This is where the Gospel's call to hear enters the life of the worshipping Church. We bring our need to the altar. The same Lord whose voice calls us is giving himself for the people he gathers.

His peace must reach the relationships we bring with us. Chrysostom presses that point upon his own congregation. They hear peace proclaimed again and again in worship. How can they receive it while deliberately keeping a feud alive? Peace with people we already like leaves the harder part of the gift unanswered.

This week, choose one place where you have been refusing to hear. It may be a fault you always explain away, or a resentment you feed by repeating the story of another person's failure. Let the King's truth name your own responsibility there. If you owe an apology, make it plain. If you have been passing on a cruel account, stop passing it on. Take one real step that lets received peace become an act.

Such a step does not mean pretending that wrongdoing never happened. Truth belongs to this kingdom. Nor can we decide another person's response. We can acknowledge the wrong we have done, resist the wrong that continues, and refuse to make hatred the principle of our conduct. Christ's authority reaches the other person too. Our work is to hear him in the responsibility that is ours.

That may leave a difficult conversation ahead. It may leave a relationship that takes a long time to mend. The Communion antiphon promises the King's blessing of peace at the end of a psalm filled with the force of God's voice amid a storm. Augustine reads its conclusion as peace in Christ while the world's troubles continue. An unsettled week does not mean that the blessing has failed.

The prayer after Communion joins the nourishment received at this altar to continuing service and the hope of reigning with Christ. We need that nourishment because hearing is a lifelong dependence. A first act of obedience does not make us self-sufficient. The King sustains the people he has called, and the kingdom celebrated today outlasts the powers that frighten us.

Return to the governor's hall. Pilate has a prisoner before him, but that prisoner's voice reaches us here. We have heard who he is: the Son through whom all things exist, the King who gives himself to make peace. We have heard what he gives: rescue, belonging, and the grace to respond. Let his truth reach the fault we have defended. Let his peace reach the relationship we have kept outside it. The King is present among his people. Today we can hear him, and follow.
